% =========================================================================
\section{Planteamiento del problema y funciones analíticas}\label{sec:planteamiento}
% =========================================================================

% -------------------------------------------------------------------------
\subsection{La ecuación}\label{sec:ecuacion}
% -------------------------------------------------------------------------
Nos limitamos a ecuaciones lineales homogéneas de segundo orden que, tras
«quitar denominadores», adoptan la forma
\begin{equation}\label{eq:edo}
  \boxed{\;P(x)\,y''(x)+A(x)\,y'(x)+B(x)\,y(x)=0\;}
\end{equation}
sobre un intervalo abierto $I\subseteq\R$, donde $P$ es un polinomio de grado
bajo (habitualmente dos, a lo sumo tres) y $A$, $B$ son funciones analíticas
en $I$ \cite{TEMAV}. Son ejemplos típicos
\[
  y''-x^{3}y'+y=0,\qquad
  y''+(\sin x)\,y=0,\qquad
  (1-x^{2})\,y''-2x\,y'+\ell(\ell+1)\,y=0 ,
\]
la última de las cuales es la \emph{ecuación de Legendre}. Para las
ecuaciones completas $P\,y''+A\,y'+B\,y=g(x)$ se aplican los métodos ya
conocidos (variación de las constantes) para obtener una solución
particular, de modo que basta estudiar la homogénea.

Los coeficientes ya no son constantes y el ansatz exponencial fracasa: en
$y''-x\,y=0$, por ejemplo, $y=\eul^{\lambda x}$ conduce a
$(\lambda^{2}-x)\,\eul^{\lambda x}=0$, que no puede cumplirse para todo $x$
con $\lambda$ constante. Necesitamos otra herramienta.

\begin{teorema}[Existencia y unicidad; recordatorio]\label{teo:eyu}
Sea $J\subseteq I$ un intervalo abierto en el que $P$ no se anula, y sea
$x_0\in J$. Para cualesquiera $\alpha,\beta\in\R$ existe una única solución
$y\in\mathcal C^{2}(J)$ de \eqref{eq:edo} con $y(x_0)=\alpha$,
$y'(x_0)=\beta$. En consecuencia, las soluciones de \eqref{eq:edo} en $J$
forman un espacio vectorial real de dimensión $2$
\textup{\cite{ELSGOLTZ}, \cite[cap.~IV]{TEMAV}}.
\end{teorema}

El \cref{teo:eyu} garantiza que las soluciones existen, pero no dice cómo
construirlas. El método de series proporcionará una construcción explícita
y, además, una propiedad cualitativa nueva: la \emph{analiticidad} de las
soluciones.

% -------------------------------------------------------------------------
\subsection{Funciones analíticas}\label{sec:analiticas}
% -------------------------------------------------------------------------

\begin{definicion}[Función analítica]\label{def:analitica}
Una función $f\colon I\to\R$ es \emph{analítica en $x_0\in I$} si existen
$R>0$ y coeficientes $(c_k)_{k\ge0}\subset\R$ tales que
\begin{equation}\label{eq:analitica}
  f(x)=\sum_{k=0}^{\infty}c_k\,(x-x_0)^k
  \qquad\text{para todo } x \text{ con } \abs{x-x_0}<R .
\end{equation}
Es decir: \textup{(i)} $f$ admite un desarrollo en serie de potencias
centrado en $x_0$ y \textup{(ii)} dicha serie \emph{converge a la propia
función} en un entorno de $x_0$. Se dice que $f$ es analítica en $I$ si lo es
en cada punto de $I$, y se escribe $f\in\Com(I)$.
\end{definicion}

La herramienta técnica que legitima todo el método es la siguiente.

\begin{teorema}[Series de potencias]\label{teo:series}
Sea $\sum_{k\ge0}c_k(x-x_0)^k$ una serie de potencias y sea $R\in[0,\infty]$
su \emph{radio de convergencia}, dado por la fórmula de Cauchy--Hadamard
\begin{equation}\label{eq:hadamard}
  \frac1R=\limsup_{k\to\infty}\abs{c_k}^{1/k}.
\end{equation}
\begin{enumerate}[label=\textup{(\roman*)}]
  \item La serie converge absolutamente si $\abs{x-x_0}<R$, y uniformemente
        en cada intervalo cerrado $\abs{x-x_0}\le r<R$; diverge si
        $\abs{x-x_0}>R$.
  \item Su suma $f$ es derivable en $(x_0-R,x_0+R)$ y se deriva
        \emph{término a término}:
        $f'(x)=\sum_{k\ge1}k\,c_k(x-x_0)^{k-1}$; la serie derivada tiene
        \emph{el mismo radio} $R$.
  \item En consecuencia $f\in\Cinf(x_0-R,x_0+R)$ y
        $c_k=f^{(k)}(x_0)/k!$: la serie es la serie de Taylor de su suma.
\end{enumerate}
\end{teorema}

\begin{proof}
(i) es el criterio de la raíz aplicado a $\sum\abs{c_k}\abs{x-x_0}^k$, junto
con el criterio M de Weierstrass. (ii) Como $k^{1/k}\to1$,
$\limsup\abs{k\,c_k}^{1/k}=\limsup\abs{c_k}^{1/k}$, de modo que la serie
derivada tiene radio $R$ y converge uniformemente en cada
$\abs{x-x_0}\le r<R$; el teorema de derivación de sucesiones de funciones
con derivadas uniformemente convergentes permite intercambiar suma y
derivada. (iii) Se obtiene iterando (ii) y evaluando en $x=x_0$.
\end{proof}

\begin{corolario}[Unicidad del desarrollo]\label{cor:unicidad}
Si $\sum_{k\ge0}c_k(x-x_0)^k=0$ para todo $\abs{x-x_0}<R$, con $R>0$, entonces
$c_k=0$ para todo $k$. En particular, si dos series de potencias coinciden en
un entorno de $x_0$, sus coeficientes son iguales uno a uno.
\end{corolario}

\begin{proof}
Por el \cref{teo:series}\,(iii), $c_k=f^{(k)}(x_0)/k!$ con $f\equiv0$.
La segunda afirmación se obtiene aplicando la primera a la diferencia.
\end{proof}

El \cref{cor:unicidad} es el principio de \emph{identificación de
coeficientes} que usaremos para obtener la relación de recurrencia. La
siguiente proposición garantiza que las operaciones algebraicas que haremos
con los coeficientes de la EDO no destruyen la analiticidad.

\begin{proposicion}[Álgebra de funciones analíticas]\label{prop:algebra}
Si $f$ y $g$ son analíticas en $x_0$, también lo son $f+g$, $f\,g$ y, si
$g(x_0)\neq0$, el cociente $f/g$. Los radios de convergencia de $f+g$ y $fg$
son al menos $\min(R_f,R_g)$.
\end{proposicion}

\begin{proof}[Comentario sobre la demostración]
Para la suma es inmediato; para el producto se usa el producto de Cauchy de
series absolutamente convergentes (\cref{lem:cauchy}). El cociente se reduce
a probar que $1/g$ es analítica cuando $g(x_0)\neq0$, lo que resulta
inmediato en variable compleja: $1/g$ es holomorfa en un disco donde $g$ no
se anula y toda función holomorfa en un disco es suma de su serie de Taylor
en él \cite[cap.~4]{AHLFORS}, \cite[§1.1]{KRANTZ}.
\end{proof}

\paragraph{Analítica es más que $\Cinf$}
La condición (ii) de la \cref{def:analitica} no es trivial: existen funciones
infinitamente derivables cuya serie de Taylor converge, pero no a la función.

\begin{proposicion}[Una función $\Cinf$ que no es analítica]\label{prop:suave}
La función
\begin{equation}\label{eq:suave}
  f(x)=\begin{cases}\eul^{-1/x^{2}}, & x\neq0,\\ 0, & x=0,\end{cases}
\end{equation}
es de clase $\Cinf(\R)$ y $f^{(k)}(0)=0$ para todo $k\ge0$. Su serie de
Taylor en $0$ es idénticamente nula y converge en todo $\R$, pero sólo
coincide con $f$ en $x=0$. Por tanto $f$ no es analítica en $0$, y
\[
  \Com(I)\subsetneq\Cinf(I).
\]
\end{proposicion}

\begin{proof}
Por inducción, para $x\neq0$ se tiene $f^{(k)}(x)=Q_k(1/x)\,\eul^{-1/x^{2}}$
con $Q_k$ un polinomio: $Q_0=1$ y
$Q_{k+1}(t)=2t^{3}Q_k(t)-t^{2}Q_k'(t)$. Como
$\lim_{x\to0}x^{-m}\eul^{-1/x^{2}}=\lim_{t\to\infty}t^{m}\eul^{-t^{2}}=0$
para todo $m\in\N$, se sigue que $f^{(k)}(x)\to0$ cuando $x\to0$. Si
$f^{(k)}(0)=0$, entonces
\[
  f^{(k+1)}(0)=\lim_{x\to0}\frac{f^{(k)}(x)-0}{x}
  =\lim_{x\to0}\frac1x\,Q_k(1/x)\,\eul^{-1/x^{2}}=0,
\]
lo que completa la inducción. La serie de Taylor es $\sum 0\cdot x^k\equiv0$,
y $f(x)>0$ para $x\neq0$.
\end{proof}

% Figura: e^{-1/x^2}, función C^infinito no analítica
\begin{figure}[ht]
\centering
\begin{tikzpicture}
  \begin{groupplot}[group style={group size=2 by 1, horizontal sep=1.6cm},
      oscuro, width=0.47\textwidth, height=5.0cm, samples=301]
    \nextgroupplot[xmin=-4, xmax=4, ymin=-0.3, ymax=1.08,
        xtick={-4,-2,0,2,4}, ytick={0,0.5,1},
        xlabel={$x$}, ylabel={$f(x)$},
        title={(a) Vista global}]
      \addplot[c1, domain=-4:4] {exp(-1/(x^2+1e-12))};
      \addplot[c3, densely dashed, domain=-4:4] {0};
      \addplot[guia, domain=-4:4, forget plot] {1};
      \node[font=\scriptsize, text=cian, anchor=west] at (axis cs:2.05,0.55) {$\eul^{-1/x^{2}}$};
      \node[font=\scriptsize, text=ambar, anchor=north] at (axis cs:0,-0.06)
        {serie de Taylor en $0$: $\equiv0$};
    \nextgroupplot[xmin=-0.5, xmax=0.5, ymin=-0.0055, ymax=0.02,
        xtick={-0.5,-0.25,0,0.25,0.5},
        ytick={0,0.01,0.02},
        yticklabel style={/pgf/number format/fixed, /pgf/number format/precision=2},
        scaled y ticks=false,
        xlabel={$x$},
        title={(b) Detalle cerca del origen}]
      \addplot[c1, domain=-0.5:0.5] {exp(-1/(x^2+1e-12))};
      \addplot[c3, densely dashed, domain=-0.5:0.5] {0};
  \end{groupplot}
\end{tikzpicture}
\caption{La función $f(x)=\eul^{-1/x^{2}}$ ($f(0)=0$) es de clase
$\Cinf$ y todas sus derivadas se anulan en el origen, de modo que su serie de
Taylor en $0$ es la función nula (trazo discontinuo). En (b) se aprecia lo
extraordinariamente plana que es $f$ cerca de $0$: en $x=\pm0{,}25$ vale
apenas $\eul^{-16}\approx1{,}1\cdot10^{-7}$.}
\label{fig:suave}
\end{figure}


\begin{observacion}[La explicación está en el plano complejo]\label{obs:suave-complejo}
Sobre el eje imaginario, $z=\iu t$, la misma expresión vale
$\eul^{-1/z^{2}}=\eul^{+1/t^{2}}$, que diverge de forma explosiva cuando
$t\to0$. Vista como función de variable compleja, $\eul^{-1/z^2}$ tiene en
$z=0$ una \emph{singularidad esencial} (\cref{sec:polos}), y ninguna serie
de potencias centrada en $0$ puede representarla. La regularidad sobre la
recta real es engañosa: la analiticidad real es una propiedad que sólo se
entiende cabalmente en $\C$.
\end{observacion}

% -------------------------------------------------------------------------
\subsection{Las soluciones se comportan mejor que los coeficientes}\label{sec:regularidad}
% -------------------------------------------------------------------------
Una idea guía de todo el tema es que \emph{integrar mejora la regularidad y
derivar la empeora}. En los puntos donde $P\neq0$, despejando en
\eqref{eq:edo},
\begin{equation}\label{eq:despeje}
  y''=-\frac{A}{P}\,y'-\frac{B}{P}\,y .
\end{equation}
Para recuperar $y$ hay que integrar dos veces, de modo que $y$ es, en un
sentido heurístico, «dos grados más regular» que los coeficientes $A/P$ y
$B/P$. El ejemplo extremo lo proporciona la delta de Dirac, que ni
siquiera es una función: en el sentido de las distribuciones,
\begin{equation}\label{eq:delta}
  y''=\delta(x)\;\Longrightarrow\;y'=\Heav(x)\ \text{(escalón)}
  \;\Longrightarrow\;y=x\,\Heav(x)\ \text{(rampa, continua)},
\end{equation}
donde se han tomado nulas las constantes de integración. Cada integración
«suaviza»: $\int\delta\,\dif x$ ya es una función acotada (con un salto), y
la siguiente integral es continua (con un pico); véase la \cref{fig:regularidad}.

% Figura: delta -> escalón -> rampa (cada integración suaviza)
\begin{figure}[ht]
\centering
\begin{tikzpicture}[x=0.95cm, y=0.95cm]
  % ---------- Panel 1: delta ----------
  \begin{scope}
    \draw[eje] (-1.7,0) -- (1.9,0) node[right, etiqgris] {$x$};
    \draw[eje] (0,-0.35) -- (0,2.4);
    \draw[-{Stealth[length=7pt,width=6pt]}, line width=1.4pt, cian] (0,0) -- (0,1.75);
    \draw[line width=1.2pt, cian] (-1.6,0) -- (1.75,0);
    \node[etiq, anchor=west] at (0.15,1.6) {$\delta(x)$};
    \node[etiqgris] at (0.05,-0.55) {$0$};
    \node[font=\small, text=gris] at (0.1,-1.25) {singular};
  \end{scope}
  \draw[flecha, line width=0.8pt, ambar] (2.25,1) -- node[above, etiqgris] {$\displaystyle\int$} (3.25,1);
  % ---------- Panel 2: escalón ----------
  \begin{scope}[xshift=5.2cm]
    \draw[eje] (-1.7,0) -- (1.9,0) node[right, etiqgris] {$x$};
    \draw[eje] (0,-0.35) -- (0,2.4);
    \draw[line width=1.4pt, magenta] (-1.6,0) -- (0,0);
    \draw[line width=1.4pt, magenta] (0,1.2) -- (1.75,1.2);
    \draw[guia] (0,0) -- (0,1.2);
    \fill[fondo] (0,0) circle (2.4pt); \draw[magenta, line width=1pt] (0,0) circle (2.4pt);
    \fill[magenta] (0,1.2) circle (2.4pt);
    \node[etiq, anchor=south] at (1.0,1.3) {$\Heav(x)$};
    \node[etiqgris, anchor=east] at (-0.08,1.2) {$1$};
    \node[etiqgris] at (0.05,-0.55) {$0$};
    \node[font=\small, text=gris] at (0.1,-1.25) {acotada, discontinua};
  \end{scope}
  \draw[flecha, line width=0.8pt, ambar] (7.45,1) -- node[above, etiqgris] {$\displaystyle\int$} (8.45,1);
  % ---------- Panel 3: rampa ----------
  \begin{scope}[xshift=10.4cm]
    \draw[eje] (-1.7,0) -- (1.9,0) node[right, etiqgris] {$x$};
    \draw[eje] (0,-0.35) -- (0,2.4);
    \draw[line width=1.4pt, ambar] (-1.6,0) -- (0,0) -- (1.6,1.6);
    \node[etiq, anchor=east] at (-0.2,1.55) {$x\,\Heav(x)$};
    \node[etiqgris] at (0.05,-0.55) {$0$};
    \node[font=\small, text=gris] at (0.1,-1.25) {continua};
  \end{scope}
\end{tikzpicture}
\caption{Cada integración gana un grado de regularidad: la delta de Dirac
(singular) se integra en el escalón de Heaviside $\Heav$ (acotado pero
discontinuo), y éste en la rampa $x\,\Heav(x)$ (continua, con derivada
discontinua). Por la \eqref{eq:despeje}, la solución $y$ se obtiene
integrando dos veces una expresión en los coeficientes.}
\label{fig:regularidad}
\end{figure}


\begin{observacion}[Singularidades fijas]\label{obs:fijas}
Para ecuaciones \emph{lineales}, las únicas singularidades posibles de las
soluciones están en los puntos donde los coeficientes de la forma normal
\eqref{eq:despeje} son singulares, es decir, en los ceros de $P$; por eso se
llaman \emph{singularidades fijas}. Las ecuaciones no lineales no tienen esta
propiedad: la solución de $y'=y^{2}$, $y(0)=1$, es $y=1/(1-x)$, que
explota en $x=1$ aunque el «coeficiente» $y^{2}$ sea perfectamente regular
\cite[cap.~3]{BENDER}. El \cref{teo:fuchs} hará precisa esta idea en el
contexto analítico.
\end{observacion}
